% !TeX root = ../main.tex
\chapter{Miscellaneous Lemmas}

\section{Many-to-one Formula}

\section{Lemmas on Log-Laplace Transforms}

    Consider the point process $\Xi$ governing the branching random walk. Let $\psi$ be the log-Laplace transform defined by
    \begin{equation}
        \psi(\lambda)
        :=
        \log\mathbb E\left[
            \sum_{x\in\Xi}e^{-\lambda x}
        \right]
        \in(-\infty,\infty],\quad
        \lambda\in\mathbb R.
    \end{equation}
    And it can be viewed as the log-Laplace transform 
    \begin{equation}
        \psi(\lambda)=\log\int_{\mathbb R}e^{-\lambda x}\,\mu(\mathrm dx)
    \end{equation}
    of the measure
    \begin{equation}
        \mu(A)
        :=
        \mathbb E\left[
            \sum_{x\in\Xi}\mathbb I[x\in A]
        \right],
        \quad
        A\in\mathscr B(\mathbb R).
    \end{equation}
    \begin{lemma} \it
        (i) The domain $\{\lambda:\psi(\lambda)<\infty\}$ is an interval and $\psi$ is convex.
        (ii) $\psi$ is $C^\infty$ on the interior of its domain.
        (iii）$\psi$ is analytic on the interior of its domain.
    \end{lemma}
    Proof is omitted. Note that (i) follows from H\"older's inequality of $\mu$; (ii) follows from Mean Value Theorem and Dominated Convergence Theorem and the fact that for any $n\ge 0$, $|x|^ne^{-\lambda x}\le C_ne^{-\lambda'x}$ for some $\lambda'>\lambda$ and constant $C_n$; and (iii) follows from Fubini-Tonelli Theorem.

    Now consider the random walk with step distribution $\xi$ obtained from the many-to-one lemma. Let $\kappa$ be the log-Laplace transform defined by
    \begin{equation}
        \kappa(\theta):=\log\mathbb Ee^{\theta\xi},\quad
        \theta\in\mathbb R.
    \end{equation}
    By (???), We can get
    \begin{equation}
        \kappa(\theta)
        =\mathbb E\left[
            \sum_{x\in\Xi}e^{\theta x}e^{-x}
        \right]
        =\psi(-\theta+1).
    \end{equation}
    We immediately obtain
    \begin{lemma}\it
        The domain of $\kappa$ is an interval, $\kappa$ is convex, and is $C^\infty$ and analytic on the interior of its domain.
    \end{lemma}

\section{Lemmas on Applying Mogul'skii's Theorem}

    Consider some independent identically distributed random variables $(\xi_i)_{i\ge 1}$ and some sequence $(x(n))$ indexed by $n$. Let
    \begin{equation}
        s_n(t):=\frac{1}{x(n)}\sum_{i=0}^{\lfloor nt\rfloor}\xi_i,\quad
        t\in[0,1],\,\xi_0=0.
        \label{eq:mogulskii-s-n-t}
    \end{equation}
    $s_n$ is thus defined as a random real function taking values in $D(0,1)$ (i.e., real functions on $[0,1]$ such that càdlàg on $[0,1)$ and left-continuous at $1$).
    
    Mogul'skii defined some functional $H_\alpha$ for $\alpha\in(0,2]$ on some set $\mathfrak N$ of $D(0,1)$. We state here only the definition for $G\in\mathfrak N$ of the form
    \begin{equation}
        G=\{f\in D(0,1):f(0)=0,\,S_2(t)<f(t)<S_1(t),\,\forall t\in[0,1]\},
        \label{eq:mogulskii-G}
    \end{equation}
    where $S_1,S_2$ are right-continuous and piecewise constant such that
    \begin{equation}
        S_2(t-0)\wedge S_2(t+0)
        <
        S_1(t-0)\wedge S_1(t+0),\quad
        t\in[0,1].
    \end{equation}
    and we have the definition
    \begin{equation}
        H_\alpha(G):=\int_0^1(S_1(t)-S_2(t))^{-\alpha}\,\mathrm dt.
        \label{eq:mogulskii-h-alpha}
    \end{equation}
    
    Recall the Mogul'skii's Theorem for $\alpha=2$. 
    
    \begin{theorem*}[Mogul'skii's Theorem]\it 
        \label{theorem:mogulskii}
        Suppose that $\mathbb E\xi_i=0$, $\mathbb E\xi_i^2=1$ and that the sequence $(x(n))$ is such that
        \begin{equation*}
            \lim_{n\to\infty}x(n)=\infty,\quad
            \lim_{n\to\infty}x(n)n^{-1/2}=0.
        \end{equation*}
        Then the relation
        \begin{equation*}
            \log\mathbb P(s_n\in G)\sim-\frac{\pi^2}{2}H_2(G)nx^{-2}(n)
        \end{equation*}
        holds for set $G\in\mathfrak N$.
    \end{theorem*}
    
    Now consider a sequence $(L_\Delta)_{\Delta\ge 0}$ indexed by $\Delta$ with $\lim_{\Delta\to\infty}L_\Delta=\infty$. Basically, when applying Mogul'skii's Theorem, we are letting $n$ there to be $L_\Delta$ and $x(n)$ there to be $\Delta$, to form an inverse sequence. We are in this section to make rigorous this statement.
    
    Let
    \begin{align}
        x(n)&:=\inf\{\Delta\ge 0: L_\Delta+1\ge n\},&n\ge 0,
        \\ y(n)&:= x(n+1)-1,&n\ge 0.
    \end{align}
    
    \begin{lemma}\it\label{lemma:mogulskii-x-y}
    The two sequences $x$ and $y$ satisfy
    \begin{equation}
        \lim_{n\to\infty}x(n)=\lim_{n\to\infty}y(n)=\infty.
    \end{equation}
    If we further require that $\lim_{\Delta\to\infty}\frac{L_\Delta}{\Delta^2}=\infty$, we have
    \begin{equation}
        \lim_{n\to\infty}x(n)n^{-1/2}
        =\lim_{n\to\infty}y(n)n^{-1/2}
        =0.
    \end{equation}
    \end{lemma}
    The proof for Lemma \ref{lemma:mogulskii-x-y} is omitted. We shall use the fact $L_{x(n)}\ge n-1$ and thus $0<\frac{x^2(n)}{n}<\frac{x^2(n)}{L_{x(n)}}$.
    
    Consider the random walk $(Y_n)_{n\ge 0}$ with zero mean and finite variance $0<\sigma^2<\infty$ obtained from the many-to-one lemma. For any $a\in(0,1)$, we denote in this section
    \begin{align}
        q(\Delta,n)&:=\log\mathbb P\left(
            -a\Delta\le Y_j\le (1-a)\Delta,\,
            \forall j\in[\![1,n]\!]
        \right),\\
        p(\Delta,n)&:=\frac{\Delta^2}nq(\Delta,n).
    \end{align}
    For any $\Delta\ge 0$, there exists $n=n(\Delta)=L_\Delta+1$ such that $x(n)\le\Delta<x(n+1)$,  and we can get by monotonicity of probability that
    \begin{lemma} \it\label{lemma:mogulskii-mono}
        For $\Delta_1\le\Delta_2$ such that $L:=L_{\Delta_1}=L_{\Delta_2}$, we have $q(\Delta_1,L)\le q(\Delta_2,L)$.
    \end{lemma}
    
    We are now able to get to the main lemma with an additional condition for $(L_\Delta)$.
    \begin{lemma}
        \it
        \label{lemma:mogulskii-cal}
        \quad
        \begin{enumerate}[(i)]
            \item By Mogul'skii's Theorem, we have, indexed by $n$,
                \begin{equation}
                    \lim_{n\to\infty}p(x(n),n-1)
                    =\lim_{n\to\infty}p(y(n),n-1)
                    =-\frac{\pi^2}2\sigma^2.
                \end{equation}
            \item For $n=n(\Delta)$, and indexed by $\Delta$,
                \begin{equation}
                    \lim_{\Delta\to\infty}p(x(n),n-1)
                    =\lim_{\Delta\to\infty}p(y(n),n-1)
                    =-\frac{\pi^2}2\sigma^2.
                \end{equation}
            \item Under the additional condition for $(L_\Delta)$ that, when $\Delta\to\infty$,
                \begin{equation}
                    N_\Delta:=\#\left\{\hat\Delta\ge 0: L_\Delta=L_{\hat\Delta}\right\}
                    =O_\Delta(1),
                \end{equation}
                we have
                \begin{equation}
                    \lim_{\Delta\to\infty}p(\Delta,L_\Delta)=-\frac{\pi^2}2\sigma^2.
                \end{equation}
            \item Under the same condition as (iii), for $v>0$,
                \begin{align}
                    &\quad
                    \lim_{\Delta\to\infty}
                    \frac{\Delta^2}{L_\Delta}
                    \log\mathbb P\left(
                        -a\Delta\le Y_j\le (1-a)\Delta,\,
                        \forall j\in[\![1,L_\Delta]\!],\,
                        Y_{L_\Delta}\in[-a\Delta,(v-a)\Delta]
                    \right)
                    \nonumber \\ &
                    =
                    -\frac{\pi^2}{2}\sigma^2.
                \end{align}
        \end{enumerate}
    \end{lemma}
    \begin{proof}[引理 \ref{lemma:mogulskii-cal}]
        For (i) and (ii), we prove here only the result for $x$. The result for $y$ is the same since the conditions for Mogul'skii's Theorem (Theorem \ref{theorem:mogulskii}) are satisfied by both $x$ and $y$ as stated in Lemma \ref{lemma:mogulskii-x-y}.
        \begin{enumerate}[(i)]
            \item  We denote here the step random variables of $(Y_i)$ as $(\sigma\xi_i)$. Define $s_n$ as in \eqref{eq:mogulskii-s-n-t}. We can express $q(x(n),n-1)$ in $s_n$ as
            \begin{align}
                q(x(n),n-1)
                &=
                \log\mathbb P
                \left(
                    \frac 1\sigma(-a)\le s_n(t)\le\frac 1\sigma(1-a),\,
                    t\in[0,1]
                \right) \nonumber
                \\ &=:\log\mathbb P(s_n\in G_0),
                \quad n\ge 1.
            \end{align}
            For each $m\ge 1$, we apply Mogul'skii's Theorem for $G=G(m)$ of the form \eqref{eq:mogulskii-G} with
            \begin{equation}
                S_1=\frac{1}\sigma\left(1-a-\frac 1{2m}\right),\,
                S_2=\frac{1}\sigma\left(-a+\frac 1{2m}\right),
            \end{equation}
            and we can use \eqref{eq:mogulskii-h-alpha} to calculate the corresponding $H_2(G)$ as
            \begin{equation}
                H_2(G)=\sigma^2\left(1-\frac 1m\right)^{-2}.
            \end{equation}
            Hence, by Mogul'skii's Theorem, we have
            \begin{align}
                \lim_{n\to\infty}
                \frac{x^2(n)}{n}
                \log\mathbb P(s_n\in G(m))
                =
                -\frac{\pi^2}2\sigma^2\left(1-\frac 1m\right)^{-2}.
            \end{align}
            Since $G(m)\subset G_0$, by monotonicity of probability, 
            \begin{align}
                \liminf_{n\to\infty}p(x(n),n-1)
                &\ge\liminf_{n\to\infty}\frac{x^2(n)}n\log\mathbb P(s_n\in G(m))
                \nonumber\\ &
                =-\frac{\pi^2}2\sigma^2\left(1-\frac 1m\right)^{-2}.
            \end{align}
            Let $m\to\infty$ and we get
            \begin{equation}
                \liminf_{n\to\infty}p(x(n),n-1)\ge-\frac{\pi^2}{2}\sigma^2.
            \end{equation}
            Similarly, we apply Mogul'skii's Theorem for $G=\widetilde G(m)$ of the form \eqref{eq:mogulskii-G} with
            \begin{equation}
                S_1=\frac{1}\sigma\left(1-a+\frac 1{2m}\right),\,
                S_2=\frac{1}\sigma\left(-a-\frac 1{2m}\right),
            \end{equation}
            and we can get
            \begin{equation}
                \limsup_{n\to\infty}p(x(n),n-1)\le-\frac{\pi^2}2\sigma^2\left(1-\frac 1m\right)^{-2}.
            \end{equation}
            Let $m\to\infty$ and the result would follow.
            \item When $\Delta\to\infty$, $n=n(\Delta)$ goes to infity because of the original condition on $(L_\Delta)$, and the result follows from (i) for subsequences.
            \item By Lemma \ref{lemma:mogulskii-mono} and since $x(n)\le\Delta\le x(n+1)-1=y(n)$, we have the lower and upper bound for $p(\Delta,L_\Delta)$ as
            \begin{align}
                p(\Delta,L_\Delta)
                &=\frac{\Delta^2}{L_\Delta}q(\Delta,L_\Delta)
                \ge\frac{\Delta^2}{L_\Delta}q(x(n),L_\Delta)
                =\frac{\Delta^2}{x^2(n)}p(x(n),n-1)
                \\
                p(\Delta,L_\Delta)
                &=\frac{\Delta^2}{L_\Delta}q(\Delta,L_\Delta)
                \le\frac{\Delta^2}{L_\Delta}q(y(n),L_\Delta)
                =\frac{\Delta^2}{y^2(n)}p(y(n),n-1)
            \end{align}
            Let $\Delta\to\infty$ in both equations and it suffices to show that
            \begin{equation}
                \lim_{\Delta\to\infty}\frac{\Delta}{x(n)}
                =\lim_{\Delta\to\infty}\frac{\Delta}{y(n)}
                =1.
            \end{equation}
            Indeed, we have $x(n)\le\Delta\le x(n)+N_\Delta$ and $y(n)-N_\Delta\le\Delta\le y(n)$, and the result follows from the additional condition.
            \item The proof is analagous to (i) to (iii). We only need to calculate the value of $H_2(G)$ for $G$ of the form \eqref{eq:mogulskii-G} with
            \begin{equation}
                S_1(t)=\begin{cases}
                    \frac 1\sigma\left(1-a\mp\frac{v\wedge 1}{2m}\right),
                    &t\in[0,\frac{n-1}n),\\
                    \frac 1\sigma\left(v-a\mp\frac{v\wedge 1}{2m}\right),
                    &t\in[\frac{n-1}n,1],
                \end{cases}
                \quad
                S_2=\frac 1\sigma\left(-a\pm\frac{v\wedge 1}{2m}\right).
            \end{equation}
            We use \eqref{eq:mogulskii-h-alpha} to calculate
            \begin{align}
                H_2(G)
                &=
                \sigma^2\left(1-\frac{v\wedge 1}m\right)^{-2}\frac{n-1}n
                +
                \sigma^2\left(v-\frac{v\wedge 1}m\right)^{-2}\frac 1n
                \\ &\sim
                \sigma^2\left(1-\frac{v\wedge 1}m\right)^{-2},
            \end{align}
            when $n\to\infty$.
        \end{enumerate}
    \end{proof}

    \begin{example}
        \label{example:mogulskii-L-1}
        For $v>0$, $u>0$ and $L_\Delta=\lfloor u\Delta^{2+v}\rfloor$, we have all the conditions satisfied, i.e.,
        \begin{equation*}
            \lim_{\Delta\to\infty}\frac{L_\Delta}{\Delta^2}=\infty,\,
            N_\Delta = O_\Delta(1).
        \end{equation*}
        Indeed, $L_\Delta>u\Delta^{2+v}-1$, and $N_{\Delta}=1$ for $\Delta$ large enough.
    \end{example}

    \begin{example}
        \label{example:mogulskii-L-2}
        For $v>0$, and $L_\Delta=\lfloor\Delta^{2+v}\rfloor-\lfloor\Delta^2\rfloor$, we have all the conditions satisfied. Indeed, $L_\Delta>\Delta^{2+v}-\Delta^2-2$, and $N_\Delta=1$ for $\Delta$ large enough.
    \end{example}